\section{Gaussian cardinal functions}\label{sec:fourier}

The construction uses a family of entire functions whose values and first
derivatives can be prescribed at a periodic set of nodes.  Multiplication by
a fixed Gaussian turns these functions into radial Schwartz functions.  Their
Fourier transforms have strictly stronger Gaussian decay, which will make the
coupled interpolation equations tractable.

Set
\begin{equation}\label{fourier:parameters}
 b=\frac{\sqrt3}{2},\qquad h=\frac25,\qquad B=b^{-2}=\frac43,
 \qquad s=b\abs{x}^2.
\end{equation}
The same coordinate, \(s=b\abs{\xi}^2\), is used on the Fourier side.
To see the lattice geometry behind this choice, let \(\Lambda\) have basis
\begin{equation}\label{fourier:lattice}
 u=\frac{1}{\sqrt b}(1,0),\qquad
 v=\frac{1}{\sqrt b}(1/2,b).
\end{equation}
This oriented basis has determinant one.  If \(\mathcal R\) is
counterclockwise rotation through \(90\) degrees, its dual basis is
\(-\mathcal R v,\mathcal R u\).  Hence
\(\Lambda^*=\mathcal R\Lambda\): the lattice and its dual have the same
shell radii.  For \(a=ju+kv\), with \(j,k\in\ZZ\), the squared radius in
the common coordinate is
\begin{equation}\label{fourier:shell-coordinate}
 s_a:=b\abs{a}^2=j^2+jk+k^2.
\end{equation}

We impose interpolation conditions at the positive nodes of a periodic set:
\begin{equation}\label{fourier:nodes}
 A=\{0,1,3,4,7,9\},\qquad
 \mathcal N=(12\ZZ+A)\cap(0,\infty).
\end{equation}
The set \(\mathcal N\) contains every positive value of \(j^2+jk+k^2\).
Indeed,
\(j^2+jk+k^2\equiv(j-k)^2\pmod3\), so it is never \(2\pmod3\).
Modulo \(4\), it is \(0\) when \(j,k\) are both even and is odd otherwise.
Combining these restrictions gives precisely the six allowed
residues in \(A\).  The inclusion is strict: \(15\in\mathcal N\), but
\(j^2+jk+k^2=15\) would imply
\((2j+k)^2\equiv2k^2\pmod5\).  Since \(2\) is not a square modulo \(5\),
this forces \(5\mid k\) and then \(5\mid j\), which would make \(25\)
divide \(15\).  Thus we require zeros at some nodes that are not lattice
shell coordinates.
This deliberate enlargement uses a periodic residue pattern with an
elementary trigonometric vanishing function. Covering planar shell values by
periodic sets also occurs in Sardari's Fourier interpolation construction
\cite[Section~6.1]{Sardari2021}; the Gaussian transformation and the
coupled inversion needed here will be proved directly.

\subsection{The cardinal family}

The values and first derivatives will be encoded by double-pole and
simple-pole terms multiplied by a common vanishing factor. This is the
same cardinal pattern that occurs in Gaussian Beurling--Selberg
interpolation \cite[formulas~(1.2)--(1.3)]{CarneiroLittmannVaaler2013}.
The nodes and simultaneous Fourier constraints here are different.

For \(-6\le j\le6\), put \(t_j=j/6\), and define
\begin{equation}\label{fourier:polynomial}
 P(s)=\prod_{a\in A}\left(2\sin\frac{\pi(s-a)}{12}\right)^2
     =\sum_{j=-6}^{6}P_j e^{i\pi t_js}.
\end{equation}
Here \(P_{-j}=\overline{P_j}\), and
\begin{equation}\label{fourier:coefficients}
 (P_0,P_1,\ldots,P_6)
 =\left(5,-1+2ib,1+2ib,-2,-\frac12+ib,-1-2ib,1\right).
\end{equation}
For example, with \(w=e^{i\pi s/6}\) and \(\rho=e^{i\pi/6}\), each factor is
\[
 \left(2\sin\frac{\pi(s-a)}{12}\right)^2
 =-\rho^{-a}w^{-1}(w-\rho^a)^2.
\]
The six minus signs cancel, and \(\rho^{-\sum_{a\in A}a}=\rho^{-24}=1\).
Thus the product in \eqref{fourier:polynomial} equals
\[
 w^{-6}\prod_{a\in A}(w-\rho^a)^2.
\]
Expanding this Laurent polynomial gives
\eqref{fourier:coefficients}.

Every zero \(n\in12\ZZ+A\) of \(P\) is double.  Define \(Q_n,D_n\) by
\begin{equation}\label{fourier:node-expansion}
 P(n+u)=Q_nu^2\bigl(1+D_nu+O(u^2)\bigr).
\end{equation}
If \(a\in A\) is the residue of \(n\) modulo \(12\), then
\begin{equation}\label{fourier:node-data}
 \begin{split}
 Q_n&=\left(\frac\pi6\right)^2
       \prod_{\substack{a'\in A\\a'\ne a}}
       \left(2\sin\frac{\pi(a-a')}{12}\right)^2,\\
 D_n&=\frac\pi6\sum_{\substack{a'\in A\\a'\ne a}}
             \cot\frac{\pi(a-a')}{12}.
 \end{split}
\end{equation}
Thus \(Q_n,D_n\) are periodic in \(n\), and their exact values are
\[
\begin{array}{c|c|c}
 a&Q_n&D_n\\\hline
 0&\pi^2/3&-7\pi\sqrt3/18\\
 1&(2\pi^2/3)(2-\sqrt3)&\pi(3+\sqrt3)/18\\
 3&(2\pi^2/3)(2-\sqrt3)&-\pi(3+\sqrt3)/18\\
 4&\pi^2/3&7\pi\sqrt3/18\\
 7&(2\pi^2/3)(2+\sqrt3)&-\pi(3-\sqrt3)/18\\
 9&(2\pi^2/3)(2+\sqrt3)&\pi(3-\sqrt3)/18
\end{array}
\]
The expansion follows by isolating the vanishing sine factor; logarithmic
differentiation of the remaining factors gives \(D_n\).  The table then uses
the elementary trigonometric values at multiples of \(\pi/12\).
In particular,
\begin{equation}\label{fourier:node-bounds}
 Q_n>\frac74=1.75,\qquad \abs{D_n}<\frac{53}{25}=2.12.
\end{equation}
To check these inequalities with rational bounds, use
\(157/50<\pi<22/7\) and \(\sqrt3<26/15\).  The minimum of \(Q_n\) is
greater than \(197192/112500>7/4\), and the maximum of \(\abs{D_n}\) is
less than \(286/135<53/25\).

Let \(C\in\RR\), and let \((c_n,d_n)_{n\in\{0\}\cup\mathcal N}\) be a
real coefficient list with
\begin{equation}\label{fourier:list-norm}
 \norm{(c,d)}_1
 :=\sum_{n\in\{0\}\cup\mathcal N}(\abs{c_n}+\abs{d_n})<\infty.
\end{equation}
Away from the nodes, define
\begin{equation}\label{fourier:cardinal}
 p(s)=P(s)R(s),\qquad
 R(s)=C+\sum_{n\in\{0\}\cup\mathcal N}
       \left(\frac{c_n}{(s-n)^2}+\frac{d_n}{s-n}\right).
\end{equation}
We next show that the apparent singularities are removable, with uniform
control of the resulting entire function.

\begin{lemma}[Finite-measure representation]\label{fourier:measure-lemma}
The function in \eqref{fourier:cardinal} extends to an entire function with
the representation
\begin{equation}\label{fourier:measure}
 p(s)=\int_{-1}^{1}e^{i\pi ts}\,\dd\mu(t),\qquad s\in\CC,
\end{equation}
where \(\mu\) is a finite complex measure, supported in \([-1,1]\), with
\(\mu(-E)=\overline{\mu(E)}\) for Borel sets \(E\).
Its total variation satisfies
\begin{equation}\label{fourier:variation}
 \norm{\mu}_{\mathrm{TV}}
 \le25\abs C+\frac{65\pi^2}{18}\sum_n\abs{c_n}
                 +\frac{32\pi}{3}\sum_n\abs{d_n}.
\end{equation}
The atoms of \(\mu\) have mass \(CP_j\) at \(t_j\), for \(-6\le j\le6\).
For real-axis evaluations, the remaining measure can be folded onto
\([0,1]\): on \(t_{j-1}<t<t_j\), its folded density is
\(\sum_n(c_nu_n(t)+d_nv_n(t))\), where
\begin{equation}\label{fourier:folded}
 \begin{split}
 u_n(t)&=-2\pi^2e^{-i\pi tn}
              \sum_{l=j}^{6}(t_l-t)P_l e^{i\pi t_ln},\\
 v_n(t)&=2i\pi e^{-i\pi tn}
              \sum_{l=j}^{6}P_l e^{i\pi t_ln}.
 \end{split}
\end{equation}
The folded atoms have masses \(CP_0\) at \(0\) and \(2CP_j\) at \(t_j\),
\(1\le j\le6\).  Integrate against this folded measure and take the real
part to obtain \(p(s)\) for real \(s\), or any of its real-axis derivatives.
\end{lemma}

\begin{proof}
Fix a node \(n\) and write \(u=s-n\).  For real \(\tau\), the identities
\begin{equation}\label{fourier:division-identities}
 \begin{split}
 \frac{e^{i\pi\tau u}-1}{u}
   &=i\pi\int_0^\tau e^{i\pi tu}\,\dd t,\\
 \frac{e^{i\pi\tau u}-1-i\pi\tau u}{u^2}
   &=-\pi^2\int_0^\tau(\tau-t)e^{i\pi tu}\,\dd t
 \end{split}
\end{equation}
hold for complex \(u\), with the values at zero taken by continuity.
When \(\tau<0\), the integrals are oriented integrals.  Inserting
\(\tau=t_l\) and multiplying by \(P_le^{i\pi t_ln}\), the subtracted
terms sum to zero because \(P(n)=P'(n)=0\).  On the positive interval
\((t_{j-1},t_j)\), only frequencies with \(l\ge j\) contribute.  The
resulting density for \(P(s)/(s-n)^2\) is \(u_n/2\), and that for
\(P(s)/(s-n)\) is \(v_n/2\).  The
negative-frequency densities are their conjugate reflections, which proves
\eqref{fourier:folded} and accounts for its factors of two.

All phase factors involving \(n\) have modulus one on the real integration
interval.  Hence the measures for the two individual basis functions satisfy
\[
 \begin{split}
 \int_0^1\abs{u_n(t)}\,\dd t
 &\le2\pi^2\sum_{l=1}^{6}\abs{P_l}
                      \int_0^{t_l}(t_l-t)\,\dd t
 =\pi^2\sum_{l=1}^{6}\abs{P_l}t_l^2
 =\frac{65\pi^2}{18},\\
 \int_0^1\abs{v_n(t)}\,\dd t
 &\le2\pi\sum_{l=1}^{6}\abs{P_l}t_l
 =\frac{32\pi}{3}.
 \end{split}
\]
These are also bounds for the total variations of the full measures on
\([-1,1]\).  The atomic measure for \(CP\) has total variation
\(\abs C\sum_{j=-6}^{6}\abs{P_j}=25\abs C\).  The coefficient
summability therefore gives convergence of the sum of basis measures in
total variation and proves \eqref{fourier:variation}.

For any compact subset of \(\CC\), the integrand in
\eqref{fourier:measure}, and each of its derivatives with respect to \(s\),
is bounded uniformly in \(t\in[-1,1]\).  The integral is consequently
entire, and the finite cardinal sums converge to it locally uniformly.
Away from the nodes, the series in \eqref{fourier:cardinal} converges
locally uniformly and has the same limit.  Thus the integral supplies its
entire extension.  Conjugate symmetry of the measure makes that extension
real on the real axis.
\end{proof}

In particular, for every integer \(k\ge0\),
\begin{equation}\label{fourier:derivative-bound}
 \abs{p^{(k)}(s)}\le\pi^k\norm{\mu}_{\mathrm{TV}}
 \qquad(s\in\RR).
\end{equation}
The finite cardinal sums and their derivatives converge locally uniformly:
the assertion for derivatives follows either directly from the measure
bound or from Cauchy's integral formula on a slightly larger compact set.
Passing their node values and derivatives to the limit in
\eqref{fourier:node-expansion} gives the cardinal identities
\begin{equation}\label{fourier:node-values}
 p(n)=Q_nc_n,\qquad p'(n)=Q_n(d_n+D_nc_n)
 \qquad(n\in\{0\}\cup\mathcal N).
\end{equation}
The other summands in \eqref{fourier:cardinal} vanish to second order at
\(n\), which explains why these identities involve only its own pair of
coefficients.

\subsection{Fourier transformation}

For a cardinal function \(p\), let
\[
 q_p(x)=e^{-\pi hs}p(s),\qquad s=b\abs{x}^2.
\]
Define, for real \(-1\le t\le1\),
\begin{equation}\label{fourier:frequencies}
 \lambda(t)=\frac{i}{b(t+ih)},\qquad
 z(t)=-\frac{B}{t+ih}-ih,
\end{equation}
and, using the measure from Lemma~\ref{fourier:measure-lemma}, put
\begin{equation}\label{fourier:kernel}
 K_p(s)=\int_{-1}^{1}\lambda(t)e^{i\pi z(t)s}\,\dd\mu(t).
\end{equation}

\begin{proposition}[Gaussian Fourier pairing]\label{fourier:transform}
The function \(q_p\) is real, radial, and Schwartz on \(\RR^2\), and
\[
 \widehat{q_p}(\xi)=e^{-\pi hs}K_p(s),\qquad s=b\abs{\xi}^2.
\]
The function \(K_p\) is entire and real on the real axis.  Moreover,
\begin{equation}\label{fourier:damping}
 \operatorname{Im}z(t)
   =h\left(\frac{B}{t^2+h^2}-1\right)
   \ge\frac{26}{435}>0\qquad(-1\le t\le1).
\end{equation}
Consequently every derivative of \(K_p(s)\) decays exponentially as
\(s\to+\infty\).
\end{proposition}

\begin{proof}
By \eqref{fourier:measure},
\[
 q_p(x)=\int_{-1}^{1}e^{-\pi b(h-it)\abs{x}^2}\,\dd\mu(t).
\]
Every spatial derivative of the integrand is a polynomial in \(x\), whose
coefficients are bounded uniformly for \(t\in[-1,1]\), times a Gaussian of
modulus \(e^{-\pi bh\abs{x}^2}\).  The finite variation of \(\mu\)
justifies differentiation under the integral and bounds every Schwartz
seminorm.  Reality follows from the real values of \(p\) on the real axis,
and radiality from the dependence on \(\abs{x}^2\).

For \(\operatorname{Re}a>0\), the two-dimensional Gaussian transform in
our Fourier convention is
\[
 \widehat{e^{-\pi a\abs{x}^2}}(\xi)
   =a^{-1}e^{-\pi\abs{\xi}^2/a}.
\]
This follows from the one-dimensional Gaussian integral by taking products
and analytically continuing in \(a\) to the right half-plane.
Apply this with \(a=b(h-it)\).  Fubini's theorem applies because the
integrands have a common integrable Gaussian majorant.  We have
\[
 a^{-1}=\lambda(t),\qquad
 -\frac{\pi\abs{\xi}^2}{a}+\pi hs=i\pi z(t)s,
\]
which gives \eqref{fourier:kernel} and the asserted Fourier transform.
The curve \(z([-1,1])\) is compact, so \(K_p\) is entire by differentiation
under its defining integral.  Also,
\[
 \lambda(-t)=\overline{\lambda(t)},\qquad
 z(-t)=-\overline{z(t)}.
\]
Thus the same folded measure and real-part convention as in
Lemma~\ref{fourier:measure-lemma} apply to \(K_p\) and its real-axis
derivatives.

Finally, the imaginary part in \eqref{fourier:damping} is smallest at
\(\abs t=1\), where its value is \(26/435\).  Differentiating
\eqref{fourier:kernel} gives, for \(s\ge0\),
\[
 \abs{K_p^{(k)}(s)}
 \le\norm{\mu}_{\mathrm{TV}}
     \max_{\abs t\le1}\bigl(\abs{\lambda(t)}\,
                              (\pi\abs{z(t)})^k\bigr)
     e^{-26\pi s/435}.
\]
This proves the decay assertion.
\end{proof}

No cardinal representation of \(K_p\) is assumed.  Only its entireness, its values and
derivatives, and the estimates obtained from \eqref{fourier:kernel} will be
used in the interpolation equations.  In particular, the measure in that
formula is supported on the real interval \([-1,1]\); later complex-disk
estimates concern specified finite integrands, not an analytic continuation
in the integration variable of an arbitrary \(\ell^1\) density.

\begin{remark}[A real-radius Hankel representation]\label{fourier:hankel}
For real \(s\ge0\), the same transform has the representation
\begin{equation}\label{fourier:hankel-formula}
 K_p(s)=e^{\pi hs}\frac{\pi}{b}\int_0^\infty
       e^{-\pi hu}p(u)
       J_0\!\left(\frac{2\pi\sqrt{us}}{b}\right)\,\dd u,
\end{equation}
where
\(J_0(v)=(2\pi)^{-1}\int_0^{2\pi}e^{-iv\cos\theta}\,\dd\theta\)
for real \(v\).  To prove this, use polar coordinates in the Fourier
integral for \(q_p\), perform the angular integral, and put \(u=b r^2\).
The resulting radial integral is
\(\widehat q_p(\xi)=(\pi/b)\int_0^\infty e^{-\pi hu}p(u)
J_0(2\pi\sqrt{us}/b)\,\dd u\), with \(s=b\abs{\xi}^2\).
Absolute convergence follows from \(\abs{p(u)}\le\norm{\mu}_{\mathrm{TV}}\)
and \(\abs{J_0(v)}\le1\).  Multiplication by \(e^{\pi hs}\) gives
\eqref{fourier:hankel-formula}.  This is an identity on the nonnegative real
axis; it is not used as a numerical sign certificate below.
\end{remark}

\subsection{The coupled interpolation conditions}

We use two cardinal functions \(p_1,p_2\), with coefficients
\(c_{i,n},d_{i,n},C_i\) for function index \(i\in\{1,2\}\).  Their
coefficients at positive nodes will be determined later; fix
\begin{equation}\label{fourier:fixed}
 (c_{1,0},d_{1,0},C_1)=(1,0.44,-0.013),\qquad
 (c_{2,0},d_{2,0},C_2)=(0,-0.368,0.017).
\end{equation}
All terminating decimals in the construction denote exact rational numbers.
Write \(K_i=K_{p_i}\) and \(H_i=p_i+K_{3-i}\).  The associated pair is
\begin{equation}\label{fourier:pair}
 F(x)=e^{-\pi hs}H_1(s),\qquad
 \widehat F(\xi)=e^{-\pi hs}H_2(s),
 \quad s=b\abs{x}^2\ \text{or}\ b\abs{\xi}^2,
\end{equation}
respectively.  Indeed, \(F=q_{p_1}+\widehat{q_{p_2}}\), and Fourier
inversion gives \(\widehat F=\widehat{q_{p_1}}+q_{p_2}\), since the
functions are even.  Thus both functions in \eqref{fourier:pair} are real,
radial, and Schwartz for every pair of admissible coefficient lists.

The required interpolation conditions are
\begin{equation}\label{fourier:interpolation}
 \begin{split}
 H_1(n)&=H_2(n)=0,\\
 H_2'(n)&=0,\qquad H_1'(n)=-Q_n\mathbf1_{\{n=1\}}
 \qquad(n\in\mathcal N).
 \end{split}
\end{equation}
Here primes denote derivatives with respect to \(s\).
They give zeros of order at least two for both functions at every positive
node except that \(H_1\) has a prescribed simple zero at \(1\).

\begin{lemma}[Coefficient form of the interpolation]\label{fourier:equivalence}
For two real \(\ell^1\) coefficient lists, conditions
\eqref{fourier:interpolation} are equivalent to
\begin{equation}\label{fourier:coefficient-equations}
 \begin{split}
 c_{i,n}&=-\frac{K_{3-i}(n)}{Q_n},\\
 d_{i,n}&=\frac{D_nK_{3-i}(n)-K_{3-i}'(n)}{Q_n}
          -\mathbf1_{\{i=1,\ n=1\}}
 \qquad(i=1,2,\ n\in\mathcal N).
 \end{split}
\end{equation}
\end{lemma}

\begin{proof}
By \eqref{fourier:node-values},
\[
 H_i(n)=Q_nc_{i,n}+K_{3-i}(n),\qquad
 H_i'(n)=Q_n(d_{i,n}+D_nc_{i,n})+K_{3-i}'(n).
\]
The first equation in \eqref{fourier:coefficient-equations} is therefore
equivalent to \(H_i(n)=0\).  Substituting it into the second displayed
identity gives the derivative condition in
\eqref{fourier:interpolation} exactly when the second coefficient equation
holds.  Division by \(Q_n\) is legitimate by
\eqref{fourier:node-bounds}.
\end{proof}
